\documentclass[12pt]{article}
\usepackage[utf8]{inputenc}
\usepackage{amsmath, amssymb, amsthm}
\usepackage{geometry}
\usepackage{fancyhdr}
\usepackage{graphicx}
\usepackage{float}

% Page geometry
\geometry{a4paper, margin=1in}

% Header and Footer setup
\pagestyle{fancy}
\fancyhf{} % Clear all header and footer fields
\lhead{Name: Nissarg Danidhariya \\ Enrollment: AU2440261}
\rhead{Course: CSE400 \\ Date: February 10, 2026}
\setlength{\headheight}{28pt} % Adjust for multiline header

\begin{document}

\begin{center}
    \textbf{\Large Lecture Scribe: Transformation of Random Variables}
\end{center}

\section*{Outline}
\begin{enumerate}
    \item Transformation of Random Variables ($Y = g(X)$)
    \item Function of Two Random Variables ($Z = X + Y$)
    \item Illustrative Examples and Derivations
\end{enumerate}

\section{Transformation of a Single Random Variable}
\subsection*{Problem Definition}
Let $X$ be a random variable (RV) where the Probability Density Function (PDF) $f_X(x)$ and Cumulative Distribution Function (CDF) $F_X(x)$ are known a priori.
We define a new random variable $Y$ through the transformation:
$$ Y = g(X) $$
\textbf{Objective:} Find the PDF of the new random variable, denoted as $f_Y(y)$.

\subsection*{Derivation: The CDF Technique}
To find $f_Y(y)$, we first determine the CDF $F_Y(y)$ and then differentiate it with respect to $y$. We consider two cases based on the monotonicity of $g(x)$.

\subsubsection*{Case 1: Monotonically Increasing Function}
If $g(x)$ is monotonically increasing, there exists a unique inverse function $x = g^{-1}(y)$.
\begin{itemize}
    \item \textbf{Step 1: Determine the CDF of Y}
    $$ F_Y(y) = \Pr(Y \le y) = \Pr(g(X) \le y) $$
    Since $g$ is increasing, the inequality direction is preserved:
    $$ F_Y(y) = \Pr(X \le g^{-1}(y)) = F_X(g^{-1}(y)) $$
    
    \item \textbf{Step 2: Differentiate to find the PDF}
    $$ f_Y(y) = \frac{d}{dy} [F_Y(y)] = \frac{d}{dy} [F_X(g^{-1}(y))] $$
    Using the chain rule:
    $$ f_Y(y) = f_X(g^{-1}(y)) \cdot \frac{d}{dy} g^{-1}(y) $$
    $$ f_Y(y) = f_X(x) \frac{dx}{dy} $$
\end{itemize}

\subsubsection*{Case 2: Monotonically Decreasing Function}
If $g(x)$ is monotonically decreasing:
\begin{itemize}
    \item \textbf{Step 1: Determine the CDF of Y}
    $$ F_Y(y) = \Pr(Y \le y) = \Pr(g(X) \le y) $$
    Since $g$ is decreasing, the inequality direction flips:
    $$ F_Y(y) = \Pr(X \ge g^{-1}(y)) = 1 - F_X(g^{-1}(y)) $$
    
    \item \textbf{Step 2: Differentiate to find the PDF}
    $$ f_Y(y) = \frac{d}{dy} [1 - F_X(g^{-1}(y))] = 0 - f_X(g^{-1}(y)) \frac{d}{dy} g^{-1}(y) $$
    $$ f_Y(y) = - f_X(x) \frac{dx}{dy} $$
\end{itemize}

\subsubsection*{General Formula}
Since the PDF must always be non-negative, we combine both cases using the absolute value:
$$ f_Y(y) = f_X(g^{-1}(y)) \left| \frac{d}{dy} g^{-1}(y) \right| $$
Or equivalently:
$$ f_Y(y) = f_X(x) \left| \frac{dx}{dy} \right| \quad \bigg|_{x = g^{-1}(y)} $$

\section{Illustrative Example}
\textbf{Problem:} Let $X$ be a random variable uniformly distributed over the interval $(-1, 1)$. Find the PDF of $Y$ given the transformation:
$$ Y = \sin\left(\frac{\pi X}{2}\right) $$

\textbf{Solution:}
\begin{enumerate}
    \item \textbf{Define $f_X(x)$:}
    $$ f_X(x) = \begin{cases} \frac{1}{2} & -1 < x < 1 \\ 0 & \text{otherwise} \end{cases} $$
    
    \item \textbf{Find the Inverse Relation:}
    $$ y = \sin\left(\frac{\pi x}{2}\right) \implies x = \frac{2}{\pi} \sin^{-1}(y) $$
    
    \item \textbf{Compute the Derivative $\frac{dx}{dy}$:}
    $$ \frac{dx}{dy} = \frac{2}{\pi} \frac{d}{dy}[\sin^{-1}(y)] = \frac{2}{\pi} \frac{1}{\sqrt{1-y^2}} $$
    
    \item \textbf{Determine Limits for Y:}
    \begin{itemize}
        \item At $x = -1$: $y = \sin(-\frac{\pi}{2}) = -1$
        \item At $x = 1$: $y = \sin(\frac{\pi}{2}) = 1$
        \item Range: $-1 < y < 1$
    \end{itemize}
    
    \item \textbf{Calculate $f_Y(y)$:}
    $$ f_Y(y) = f_X(x) \cdot \left| \frac{dx}{dy} \right| $$
    $$ f_Y(y) = \frac{1}{2} \cdot \left| \frac{2}{\pi \sqrt{1-y^2}} \right| $$
    $$ f_Y(y) = \frac{1}{\pi \sqrt{1-y^2}}, \quad -1 < y < 1 $$
\end{enumerate}

\section{Function of Two Random Variables}
\textbf{Context:} We consider a derived random variable $Z$ which is a function of two random variables $X$ and $Y$:
$$ Z = g(X, Y) $$
Specifically, we analyze the sum:
$$ Z = X + Y $$

\subsection*{Derivation of the CDF}
To find the distribution of $Z$, we start with the CDF formulation:
$$ F_Z(z) = \Pr(Z \le z) = \Pr(X + Y \le z) $$

This probability is the double integral of the joint PDF $f_{XY}(x, y)$ over the region where $x + y \le z$ (equivalently $y \le z - x$).

\textbf{Integral Setup:}
$$ F_Z(z) = \int_{x=-\infty}^{\infty} \int_{y=-\infty}^{z-x} f_{XY}(x, y) \, dy \, dx $$

\section{Class Problems (Exam Focus: 10 Marks)}
The following problems were presented for solution and proofs:

\begin{enumerate}
    \item \textbf{Sum of Independent Normal RVs:}
    Let $X \sim \mathcal{N}(0, 1)$ and $Y \sim \mathcal{N}(0, 1)$ be independent random variables.
    \textbf{Prove that:} 
    $$ Z = X + Y \sim \mathcal{N}(0, 2) $$
    
    \item \textbf{Sum of Exponential RVs:}
    If $X$ and $Y$ are independent exponential random variables with parameter $\lambda$, find the PDF $f_Z(z)$.
    
    \item \textbf{General Independence Case:}
    Find $f_Z(z)$ if $X$ and $Y$ are independent. (This generally leads to the convolution of the marginal PDFs).
\end{enumerate}

\end{document}